% ============================================================
% Abstract — English
% ============================================================

\chapter*{Abstract}
\addcontentsline{toc}{chapter}{Abstract}

Agricultural monitoring in Saharan environments presents challenges that distinguish it
from conventional farming contexts: intermittent electrical supply, unreliable Internet
access, extreme thermal conditions, and the practical difficulty of reaching dispersed
field plots between observation cycles. For alfalfa cultivation specifically, these
constraints are compounded by the absence of IoT monitoring systems designed for this
crop in Saharan agricultural settings and by a recurring limitation identified across
reviewed works: the conflation of the time of data upload with the time of
measurement, which undermines the traceability of any time-series analysis.

This thesis presents the design and implementation of an IoT-based monitoring and
deterministic agronomic analytics platform for alfalfa cultivation in the El~Oued
region of Algeria. The system consists of three embedded sensor nodes communicating
over a 433~MHz LoRa wireless link within a structured ten-minute superframe. The MAIN
gateway node coordinates reception, maintains an authoritative DS3231 real-time clock,
and persists all measurements locally on an SD card independent of network
availability. Node2 acquires soil moisture, temperature, and electrical conductivity
at a second irrigation pivot via RS485/Modbus; Node3 acquires atmospheric temperature,
humidity, and pressure. Data transfer follows an offline-first model: accumulated SD
records are uploaded in batch to a cloud backend when connectivity is intermittently
available. Eight firmware operating modes address the specific constraints of the
deployment, including a recovery mechanism that corrects timing drift causing
persistent communication failure, and power-saving strategies that extend battery
autonomy through deep sleep and relay-controlled sensor power.

The backend, built with Node.js and PostgreSQL, enforces a strict policy separating
the RTC-derived measurement timestamp (\texttt{measured\_at}) from the upload receipt
time, which is stored exclusively as transfer audit metadata. A React dashboard
presents sensor trends, agronomic events, and the output of a ten-stage deterministic
quality-control and reliability-scoring pipeline. AI-assisted interpretation is
provided through a backend-proxied Gemini interface subject to six architectural
constraints that prevent the model from computing primary analytics or modifying stored
data.

A pilot validation window spanning 2026-05-06 to 2026-05-11 produced 2433 sensor
readings distributed equally across the three nodes, six completed upload sessions,
and 135 rows of deterministic pipeline output. These results demonstrate the functional
feasibility of the end-to-end architecture within the scope of a pilot deployment;
full-season agronomic validation and reliability threshold calibration constitute
priorities for future work.

\vspace{1em}
\noindent\textbf{Keywords:} Internet of Things, smart agriculture, alfalfa, LoRa,
ESP32, offline-first, deterministic analytics, Saharan agriculture, data quality,
real-time clock, timestamp traceability.
